\chapter{Nachfolger, Vorgänger und Induktion}
\label{b10guide:ch:induction}

\section{Die Gestalt natürlicher Zahlen}

\begin{theorem}[Null oder eindeutiger Nachfolger]
\label{b10guide:thm:zero-or-successor}
Für jedes $n\in\mathbb N$ gilt genau eine der beiden Möglichkeiten:
\[
 n=0
 \qquad\text{oder}\qquad
 \exists!m\in\mathbb N\;(n=S(m)).
\]
Insbesondere ist $S[\mathbb N]=\mathbb N_+$.
\end{theorem}

\begin{theorem}[Nachfolgervergleich und Ausschluss von Fixpunkten]
\label{b10guide:thm:successor-comparison}
Für alle $m,n\in\mathbb N$ gilt
\[
 S(m)=S(n)\ \Longleftrightarrow\ m=n,
 \qquad
 S(m)\ne S(n)\ \Longleftrightarrow\ m\ne n.
\]
Außerdem gilt $S(n)\ne n$ für jedes $n\in\mathbb N$.
\end{theorem}

\begin{definition}[Vorgänger einer positiven natürlichen Zahl]
\label{b10guide:def:predecessor}
Für $n\in\mathbb N_+$ bezeichnet $\operatorname{pred}(n)$ die eindeutig
bestimmte natürliche Zahl $m$ mit $S(m)=n$. Damit ist
$\operatorname{pred}\colon\mathbb N_+\to\mathbb N$;
ein Wert $\operatorname{pred}(0)$ wird nicht definiert.
\end{definition}

\begin{theorem}[Kennzeichnende Vorgängergleichungen]
\label{b10guide:thm:predecessor}
Für $m\in\mathbb N$ und $n\in\mathbb N_+$ gilt
\[
 \begin{gathered}
 \operatorname{pred}(S(m))=m,\qquad S(\operatorname{pred}(n))=n,\\
 \operatorname{pred}(n)=m\ \Longleftrightarrow\ n=S(m).
 \end{gathered}
\]
\end{theorem}

\section{Induktionsformen}

Es gilt die natürliche Ordnung aus Definition~\ref{b10guide:def:natural-order}.
$P(n)$ bezeichnet eine Eigenschaft, gegebenenfalls mit festen Parametern;
formal kann sie als $n\in U$ mit $U\subseteq\mathbb N$ gelesen werden.

\begin{theorem}[Vollständige Induktion]
\label{b10guide:thm:induction}
Gelten
\[
 P(0),\qquad
 \forall n\in\mathbb N\;(P(n)\Rightarrow P(S(n))),
\]
so gilt $P(n)$ für jedes $n\in\mathbb N$.
\end{theorem}

\begin{theorem}[Induktion ab einem Anfangswert]
\label{b10guide:thm:tail-induction}
Sei $a\in\mathbb N$. Gelten $P(a)$ und
\[
 \forall n\in\mathbb N\;
 \bigl((a\le n\land P(n))\Rightarrow P(S(n))\bigr),
\]
so gilt $P(n)$ für jede natürliche Zahl $n\ge a$.
\end{theorem}

\begin{theorem}[Starke Induktion]
\label{b10guide:thm:strong-induction}
Gilt
\[
 \forall n\in\mathbb N\;
 \Bigl(\bigl[\forall m\in\mathbb N\;(m<n\Rightarrow P(m))\bigr]
      \Rightarrow P(n)\Bigr),
\]
so gilt $P(n)$ für jedes $n\in\mathbb N$.
Für $n=0$ ist die in eckigen Klammern stehende Bedingung leer;
der Induktionsschritt verlangt an dieser Stelle bereits $P(0)$.
\end{theorem}

\begin{theorem}[Induktion auf einem natürlichen Intervall]
\label{b10guide:thm:interval-induction}
Seien $a,b\in\mathbb N$ mit $a\le b$. Gelten $P(a)$ und
\[
 \forall n\in\mathbb N\;
 \bigl((a\le n<b\land P(n))\Rightarrow P(S(n))\bigr),
\]
so gilt $P(n)$ für alle $n\in\mathbb N$ mit $a\le n\le b$.
\end{theorem}

\section{Kleinstes Element und Abstieg}

\begin{theorem}[Kleinste natürliche Zahl mit einer Eigenschaft]
\label{b10guide:thm:least-number}
Jede nichtleere Teilmenge $U\subseteq\mathbb N$ besitzt genau ein
kleinstes Element, also genau ein $u\in U$ mit
\[
 \forall v\in U\;(u\le v).
\]
Insbesondere gibt es zu einer Eigenschaft, die nicht für alle
natürlichen Zahlen gilt, eine kleinste natürliche Zahl, für die sie
nicht gilt.
\end{theorem}

\begin{theorem}[Prinzip des unendlichen Abstiegs]
\label{b10guide:thm:descent}
Sei $U\subseteq\mathbb N$. Gilt
\[
 \forall n\in U\;\exists m\in U\;(m<n),
\]
so ist $U=\varnothing$.
\end{theorem}
